\documentclass[11pt]{article}
\usepackage[margin=2.5cm]{geometry}
\usepackage{xcolor}
\usepackage{enumitem}
\usepackage{hyperref}
\usepackage{booktabs}

\definecolor{revcolor}{RGB}{0,0,160}
\definecolor{respcolor}{RGB}{0,100,0}

\newcommand{\RC}[1]{\noindent\textcolor{revcolor}{\textbf{Reviewer Comment:} \emph{#1}}\par\medskip}
\newcommand{\AR}[1]{\noindent\textcolor{respcolor}{\textbf{Author Response:} #1}\par\medskip}
\newcommand{\CH}[1]{\noindent\textbf{Changes Made:} #1\par\bigskip}

\title{Response to Reviewers\\[4pt]
\large Manuscript \#Access-2025-52745\\
``Topology-Constrained NeuroB-Rep: Contact-Aware CAD\\
Reconstruction from 3-D Point Clouds''}
\author{}
\date{}

\begin{document}
\maketitle

We sincerely thank both reviewers for their thorough and constructive feedback.
Every comment has been addressed in the revised manuscript.
Below we respond to each point, describe the specific changes made, and
provide cross-references to the relevant sections, tables, and figures
in the \textbf{revised} manuscript.

\bigskip\hrule\bigskip

%% ================================================================
\section*{Reviewer 1}
%% ================================================================

% ---- R1-M1 ----
\subsection*{R1-M1: Clarity of Contribution vs.\ Point2CAD}

\RC{The distinction between the proposed method and Point2CAD is not
sufficiently clear. A side-by-side comparison table or flowchart would
help readers understand the novelty.}

\AR{We have added Table~2 (Pipeline Comparison), which contrasts
NeuroB-Rep and Point2CAD across seven design dimensions (topology
handling, contact objective, gating schedule, continuity enforcement,
freeform surfaces, export format, and evaluation protocol).
The Introduction (Sec.~I) has been restructured so that the
industrial motivation paragraph precedes the contribution list,
making the positioning against Point2CAD explicit from the outset.}

\CH{New Table~2 in Sec.~IV; Introduction restructured (Sec.~I,
paragraphs 3--5).}

% ---- R1-M2 ----
\subsection*{R1-M2: Evaluation Protocol Consistency}

\RC{The evaluation protocol mixes different scales and metrics.
A rationale for this hybrid approach is needed.}

\AR{We have added a dedicated ``Rationale for Hybrid Evaluation''
paragraph in Sec.~V-B explaining why contact-gap energies are
evaluated in a normalized scale while geometric distances are reported
in millimeters: contact tolerances are scale-invariant design intents,
whereas Chamfer/Hausdorff must be in physical units to be
interpretable by manufacturing engineers. We also clarify the
ICP alignment step and its role in bridging the two scales.}

\CH{New paragraph ``Rationale for Hybrid Evaluation'' in Sec.~V-B.}

% ---- R1-M3 ----
\subsection*{R1-M3: Missing Baseline Comparison}

\RC{Table~3 reports only your method. At least one baseline (e.g.,
Point2CAD) should be included under the same protocol.}

\AR{Table~7 (Geometry Consistency after ICP) now includes a
Point2CAD column. We ran Point2CAD on the same three components
(Impeller, Shaft, Casing) under the identical ICP protocol
with 120 independent repeats per component.
NeuroB-Rep consistently outperforms Point2CAD: geometric accuracy
and surface coincidence errors are 2--3$\times$ lower, while length
errors show 4--8$\times$ improvement.
(Note: The table was Table~3 in the original submission; it is now
Table~7 due to newly added tables.)}

\CH{Table~7 extended with Point2CAD column (Sec.~VI-G).}

% ---- R1-M4 ----
\subsection*{R1-M4: Failure Cases}

\RC{The paper does not discuss failure modes. Please analyze cases
where the method does not perform well.}

\AR{We have added Sec.~VIII-E ``Failure Modes and Mitigations''
analyzing the cross-boundary misassignment observed in
Figs.~5--6. The primary failure originates in the upstream
PointNet segmentation stage when geometrically similar adjacent
surfaces receive the same label. We describe the downstream
effects (degraded data-term fit, missing contact edges) and
propose three mitigations: adaptive layer thresholds,
boundary-aware resampling, and multi-scale segmentation.}

\CH{New Sec.~VIII-E with failure analysis and three mitigation
strategies.}

% ---- R1-m1 ----
\subsection*{R1-m1: Abstract Wording}

\RC{``0.01\,mm'' in the abstract should reflect the range across
components.}

\AR{Changed to ``0.01--0.03\,mm'' in the abstract and corresponding
results text.}

\CH{Abstract and Sec.~VI updated.}

% ---- R1-m2 ----
\subsection*{R1-m2: Formula Notation}

\RC{The relationship between $d_{ij,k}$ and $s_{ij,k}$ after Eq.~(2)
is unclear.}

\AR{Added a clarifying sentence after Eq.~(2) linking
$d_{ij,k}$ (gap distance) and $s_{ij,k}$ (projected overlap
distance) and explaining their complementary roles in the
contact objective.}

\CH{Sentence added after Eq.~(2) in Sec.~IV-B.}

% ---- R1-m3 ----
\subsection*{R1-m3: Figure Reference Order}

\RC{Figures are not referenced in sequential order.}

\AR{All figure references have been checked and corrected to
follow sequential order throughout the manuscript.}

\CH{Cross-references updated throughout.}

\bigskip\hrule\bigskip

%% ================================================================
\section*{Reviewer 2}
%% ================================================================

% ---- R2-C1 ----
\subsection*{R2-C1: Comparison with Six Cited Methods}

\RC{The paper should compare with Li et al.\ 2019,
ComplexGen, PrimitiveNet, Supervised Fitting, ParSeNet,
and BrepGen, not just cite them.}

\AR{We provide two levels of comparison:
(1)~Table~1 (Feature Comparison) gives a qualitative
comparison across all surveyed methods on seven capability
axes (topology support, freeform surfaces, STEP export, etc.).
(2)~Table~7 provides a full quantitative comparison with
Point2CAD under identical protocol.
A new Sec.~VIII-B (``Scope of Baseline Comparisons'')
explains why the remaining six methods---Supervised
Fitting, PrimitiveNet, ParSeNet, ComplexGen, BrepGen,
and Maurell et al.---cannot be compared
numerically: dataset mismatch (ABC/ShapeNet/synthetic vs.\
industrial scans), output-format disparity (chain complexes,
latent codes vs.\ STEP), and evaluation-metric incompatibility
(IoU/F-score vs.\ mm-locked Chamfer/Hausdorff + contact-gap).
Point2CAD is the only candidate sharing the same pipeline
architecture, publicly available code, and compatible input format.}

\CH{New Table~1 (Sec.~III); extended Table~7 (Sec.~VI-G);
new Sec.~VIII-B explaining comparison scope.}

% ---- R2-C2 ----
\subsection*{R2-C2: Paper Outline Missing}

\RC{Add a paper outline at the end of the Introduction.}

\AR{A paper-outline paragraph has been added at the end of Sec.~I,
listing all sections and their content.}

\CH{Outline paragraph added at end of Sec.~I.}

% ---- R2-C3 ----
\subsection*{R2-C3: Motivation Paragraph Position}

\RC{The industrial motivation should appear before the contribution
list, not after.}

\AR{The motivation paragraph has been moved to precede the
contribution list in Sec.~I.}

\CH{Paragraph reordered in Sec.~I.}

% ---- R2-C4 / R2-C5 ----
\subsection*{R2-C4 \& R2-C5: Related Work Structure and Depth}

\RC{The Related Work section has a single subsection and lacks
detailed positioning against existing methods.}

\AR{Sec.~III (Related Work) has been restructured into four
subsections: (A)~Point-Cloud Segmentation,
(B)~Primitive Fitting and Assembly, (C)~B-Rep Learning and
Generative Models, and (D)~Hybrid Reconstruction and Point2CAD.
Each subsection discusses the strengths and limitations of prior
methods and explicitly positions our approach.
Table~1 provides a structured feature comparison.
New references have been added: Li et al.\ 2019, PrimitiveNet
(Huang 2021), Maurell et al.\ 2024.}

\CH{Sec.~III restructured with four subsections and Table~1.}

% ---- R2-C6 ----
\subsection*{R2-C6: Face/Edge/Vertex Recovery from Points}

\RC{How are faces, edges, and vertices recovered from the
unstructured point cloud?}

\AR{Sec.~IV-A now includes a dedicated paragraph
(``Patch Graph Construction'') explaining the kNN co-label adjacency
mechanism: for each point, the $k$~nearest neighbors are queried;
if a neighbor carries a different patch label, an edge is recorded
between the two patches. This yields the face adjacency graph
whose edges define candidate contact boundaries.}

\CH{New paragraph in Sec.~IV-A-2 (``From Points to a Patch Graph'').}

% ---- R2-C7 ----
\subsection*{R2-C7: Single Subsection in Method}

\RC{Section~III-A contains only one subsection, which violates
standard structuring conventions.}

\AR{The Method section (now Sec.~IV after the insertion of a new
Preliminaries section) has been expanded with multiple subsections
(IV-A through IV-F), resolving the single-subsection issue.
The original Related Work (now Sec.~III) was similarly restructured
into four subsections (III-A through III-D).}

\CH{Sec.~IV-A now contains two subsubsections (IV-A-1 and IV-A-2),
resolving the single-subsection issue; Sec.~III (Related Work)
restructured into four subsections (A--D).}

% ---- R2-C8 ----
\subsection*{R2-C8: Jang et al.\ Explanation}

\RC{The description of Jang et al.\ [12] is too brief.
Expand the explanation of the mesh repair principles.}

\AR{Sec.~II-C now provides a three-step description of Jang et
al.'s method: (i)~detection of interpenetrating triangles via
oriented bounding-box pre-filtering, (ii)~computation of exact
intersection curves between conflicting triangles, and
(iii)~local re-meshing by splitting edges along intersection
curves and re-triangulating affected patches. We note the method
operates in a single pass without global re-meshing, making it
suitable as an online diagnostic.}

\CH{Sec.~II-C expanded from one sentence to a full paragraph.}

% ---- R2-C9 ----
\subsection*{R2-C9: Projection Operator Clarity}

\RC{The passage describing the projection operator is hard to follow.}

\AR{The projection operator description has been rewritten for
clarity, with explicit notation for the nearest-point projection
onto each parametric surface and the role of the projection in
computing gap distances.}

\CH{Projection passage rewritten in Sec.~IV-B.}

% ---- R2-C10 ----
\subsection*{R2-C10: ``Soft Barrier'' Definition}

\RC{Define ``soft barrier'' at first use.}

\AR{The term is now defined in Sec.~II-C (Preliminaries) at its
first occurrence: ``a smooth penalty function
$\log(1+\exp(\tau-s))$ that grows rapidly as the distance~$s$
between projected points on two patches decreases below
threshold~$\tau$.''}

\CH{Definition added in Sec.~II-C.}

% ---- R2-C11 ----
\subsection*{R2-C11: Inconsistent Paragraph Labels}

\RC{The ``a:'' style paragraph labels are inconsistent.}

\AR{All paragraph labels have been unified using the standard
\texttt{\textbackslash paragraph\{\}} command throughout the
manuscript.}

\CH{Paragraph labels standardized throughout.}

% ---- R2-C12 ----
\subsection*{R2-C12: Fig.~3 Panel Count}

\RC{The caption references panels (a,b,c,d) but only three
panels are visible.}

\AR{This was a caption error. The original figure has always
contained three panels. The caption has been corrected to
reference (a), (b), and (c) only.}

\CH{Fig.~3 caption corrected to three panels.}

% ---- R2-C13 ----
\subsection*{R2-C13: PCA Normalization Justification}

\RC{The PCA canonical transform is used without justification.
An ablation would strengthen the claim.}

\AR{We conducted a PCA ablation experiment (5~seeds, 60~iterations
each) comparing convergence with and without PCA normalization.
Results are reported in Table~6 (Sec.~VI-B). PCA normalization
yields a modest but consistent improvement: Chamfer distance
improves by 0.3\% and Hausdorff by 1.4\%, with lower variance
across seeds. The primary benefit is stabilizing the initial
fit and reducing sensitivity to scan orientation.}

\CH{New Table~6 and Sec.~VI-B ``Effect of PCA Canonical Transform.''}

% ---- R2-C14 ----
\subsection*{R2-C14: ISO~5459 Datum Explanation}

\RC{The reference to ISO~5459 needs more context.}

\AR{Sec.~V-A now includes an expanded explanation of the
ISO~5459 datum framework: how datum features are established
from physical surfaces, and how this relates to our
evaluation's reference geometry and coordinate alignment via ICP.}

\CH{Expanded explanation in Sec.~V-A.}

% ---- R2-C15 ----
\subsection*{R2-C15: GPU Implementation Details}

\RC{No information is provided about the GPU implementation.}

\AR{We have added implementation details in Sec.~V-E
(Computational Cost) and Table~4: all optimizations run on a
single NVIDIA RTX~3090 using PyTorch autograd with CUDA-batched
energy evaluations. NeuroB-Rep requires ${\sim}7$\,min / 4.8\,GB
per component vs.\ Point2CAD's ${\sim}4$\,min / 3.2\,GB.
The overhead stems from the contact-gap and continuity terms,
which require pairwise boundary-point queries absent in
Point2CAD's distance-only objective.}

\CH{New Sec.~V-E and Table~4.}

% ---- R2-C16 ----
\subsection*{R2-C16: Artifact URL}

\RC{Provide a public artifact URL for reproducibility.}

\AR{The solver code, evaluation scripts, and default configuration
files are now publicly available at
\url{https://github.com/JuneWooKang/TopoConstrained-NeuroBrep}.
Due to proprietary constraints, the industrial scan data cannot
be included; however, all scripts reproduce the reported metrics
when provided with compatible point-cloud inputs.}

\CH{URL added in Sec.~VII.}

% ---- R2-C17 ----
\subsection*{R2-C17: Table~3 Baseline Column}

\RC{Table~3 should include at least one baseline for comparison.}

\AR{The consistency table (now Table~7 in the revised manuscript)
includes a Point2CAD column evaluated under the identical ICP
protocol (120 repeats per component).
See R1-M3 above for details.}

\CH{Point2CAD column added to Table~7 (Sec.~VI-G).}

% ---- R2-C18 ----
\subsection*{R2-C18: ``Reproducibility Hygiene'' Wording}

\RC{The term ``reproducibility hygiene'' is informal.}

\AR{Changed to ``reproducibility protocol'' throughout the
manuscript.}

\CH{Terminology updated in Sec.~VII.}

\bigskip\hrule\bigskip

%% ================================================================
\section*{Additional Corrections}
%% ================================================================

During revision, we identified and corrected the following issues:

\begin{itemize}
  \item \textbf{Iteration count (1{,}000 $\rightarrow$ 120):}
        The original submission stated ``1{,}000 independent repeats''
        in the evaluation description. This was a typographical error.
        All experiments were conducted with 120 independent repeats
        per component, as correctly reflected in the Table~7 caption.
        The text has been corrected throughout.

  \item \textbf{Fig.~3 caption (4 $\rightarrow$ 3 panels):}
        The original image has always contained three panels.
        The caption erroneously referenced four panels (a,b,c,d).
        The caption now correctly references panels (a), (b), (c).

  \item \textbf{Abstract grammar:}
        ``a based on temperature gating mechanism'' corrected to
        ``a temperature-based gating mechanism.''

  \item \textbf{Computational cost reporting:}
        Added Table~4 reporting training time, GPU memory, and
        per-component overhead for both NeuroB-Rep and Point2CAD
        on an NVIDIA RTX~3090.

  \item \textbf{New Acknowledgment section:}
        Added NRF funding acknowledgment before the References.

  \item \textbf{Section renumbering:}
        A new Sec.~II (Preliminaries) was inserted, shifting all
        subsequent section numbers by one. Correspondingly, four
        new tables (Tables~1, 2, 4, 6) were added, renumbering
        existing tables.
\end{itemize}

\bigskip
\noindent We believe these revisions comprehensively address all
reviewer concerns. We are grateful for the opportunity to
strengthen the manuscript and welcome any further suggestions.

\end{document}
